\section{Relationship to Existing Protocols}
\label{sec:mapping}
% ================================================================================

This section answers RQ1--RQ4 by comparing each Bridge Contract
obligation with the cited drafts and one product, and states the
complementarity boundary explicitly. Section~\ref{sec:related-work} gives
the per-draft readings; this section holds the cross-draft comparison and
does not restate them.

\subsection{Method}\label{sec:comparison-method}

The candidates are the drafts that name a session, resumption, human
escalation or audit concern close to the contract (AGTP with its session
and bindings companions, GAR, HEM, PACE, the Agent Audit Trail) and one
product that implements detach (Herdr). The drafts on evidence binding and
attestation composition were read and left out because they concern none of
the obligations. AgentThread analyses protocols but specifies none, so it
has no column. Termenbox, the author's own product, is closed source: no text
of it can be checked by a reader, so it has no column and claims about it stay
as reported by the author.

The rows are the contract's own obligations. Each cell records one status:
specified, partly (a neighbouring mechanism covers part),
described (a product behaviour with no specification), not specified, or
not verified. ``Not specified'' means \emph{as read}: the saved text
under \texttt{paper/sources/} was searched for the concern and nothing was
found; it is a statement about those bytes and not about the protocol's
intent or any later revision. ``Not verified'' marks a reading whose
normative text was not read, and is a stated limit rather than a guess. The
cells were first drafted on 28 October 2026; each was checked again
against the saved text on 10 November 2026, with one exception: the RESUME
method of the AGTP standard-methods companion was not read (its bytes were not
retrievable, see \texttt{paper/sources/MANIFEST.md}) and is cited as reported. The data (Table~\ref{tab:mapping}), with the version, section,
evidence and reading date of each cell, are in
\texttt{data/related-work-matrix.json}; the table is generated from it
(\texttt{make related-matrix}) and \texttt{make validate} fails on a stale
file, a cell without a section pointer or a source that is not saved.

\subsection{Matrix}

% Generated by scripts/related_work_matrix.py from data/related-work-matrix.json. Do not edit by hand.
\begin{table}[htbp]
\centering
\caption{Bridge Contract obligations against the cited drafts and Herdr, as read. S: a mechanism is specified; P: a neighbouring mechanism covers part; D: a product describes the behaviour, with no specification; $-$: not specified as read (the saved text was searched and nothing found); ?: not verified. Pointers are sections of the draft in the column; the version read, the evidence and the reading date of each cell are in \texttt{data/related-work-matrix.json}.}
\label{tab:mapping}
\scriptsize
\begin{tabularx}{\columnwidth}{p{0.17\columnwidth}>{\raggedright\arraybackslash}X>{\raggedright\arraybackslash}X>{\raggedright\arraybackslash}X>{\raggedright\arraybackslash}X>{\raggedright\arraybackslash}X>{\raggedright\arraybackslash}X}
\toprule
\textbf{Obligation} & \textbf{AGTP} & \textbf{GAR} & \textbf{HEM} & \textbf{PACE} & \textbf{Audit Trail} & \textbf{Herdr} \\
\midrule
Detached execution & P~SESSION \S3.2 & $-$ & $-$ & P~\S4.5, \S10.1 & $-$ & D~README, state page \\
Human-pending queue & P~\S8.2.10 & $-$ & P~\S4, \S5, \S13.1 & $-$ & $-$ & P~README \\
Resume proof & P~\S8.2.12 & $-$ & $-$ & P~\S4.5, Table 17 & P~\S7.7 & $-$ \\
Audit continuity & P~\S6.5.4 & S~\S3, \S6 & P~\S3.4.3 & $-$ & S~\S4, \S5.2, \S6, \S8.2 & $-$ \\
Authority drift & P~\S6.5.3, \S6.6 & P~\S2, \S12.1 & ? & P~\S9.1.2 & $-$ & $-$ \\
Detach and human timers & P~\S8.2.12 & $-$ & S~\S14.1, \S14.2 & P~\S4.6, \S10.1 & P~\S7.4 & $-$ \\
Answer binding & P~\S8.2.11 & $-$ & S~\S12.8 & $-$ & P~\S7.3 & $-$ \\
Quarantine & $-$ & $-$ & $-$ & $-$ & $-$ & $-$ \\
Restart recovery & P~SESSION \S3.2 & $-$ & $-$ & P~\S4.5, \S10.1 & P~\S8.2, \S8.3 & D~state page \\
\bottomrule
\end{tabularx}
\end{table}


\subsection{Analysis}

\emph{Neighbours, not owners.} For detached execution, the human-pending
queue, resume proof and restart recovery no source specifies
the obligation for a bridge: the best cell is partly, and the closest
mechanisms are AGTP SUSPEND with its nonce, PACE re-subscription, and the
Agent Audit Trail's pause, resume and recovery records. For quarantine
every cell is not specified: GAR quarantines telemetry spans and the Audit
Trail flags a broken chain, and neither isolates a session. Each is
a message or a record, not a check made against a durable chain, which is
why the contract attaches the proof to the log (I4, I5, I11, I12).

\emph{Overlaps the contract reuses.} Audit continuity is specified by
GAR and the Agent Audit Trail, and the human timers and the single-use
signed decision by HEM. These are the contract's I11, I14 and timer rules in
another setting, and the contract does not present them as new; it borrows
the chain shape and ties the timers to a detach. The GAR and Audit Trail
cells also conflict with one another: GAR presupposes a GEC as sole source,
while the Audit Trail lets the agent record itself and asks, from trust
level L2, for an independent recorder (an API gateway or reverse proxy, a
sidecar or middleware, or an external monitor) that SHOULD sign records
with its own key (Audit Trail \S 5.2); the bridge is a recorder of that
kind, so the Audit Trail neighbours the contract's audit role more closely
than GAR does. Their chains use different parent
fields; no mapping between them is specified, and the contract's
sequence-numbered events fit neither exactly.

\emph{Gaps.} No source checks that authority did not grow
while the client was away. AGTP does reject a request outside the
Authority-Scope with status 262 (AGTP \S 6.5.3, \S 6.6, as read); the claim is that
the check is per request and has no relation to a detach. HEM's cell for
drift is not verified. The answer-binding cell for PACE and Herdr and the timer cell for
Herdr are empty only as far as the pages read go.

On the drafts as read, every obligation has a neighbour in at least one source and
none specifies the obligation for a bridge; this is
a reading, not a proof, and it does not hold for every future revision. The
matrix has limits of its own. The obligations are the contract's, so a row
can show a neighbour or a gap but not that the contract chose the right
rows, and a cell can be wrong in either direction. The Bridge Contract is
thus positioned as a complement, not a competitor. RQ1's answer: yes,
uncovered bridge obligations exist (the gap classes G1--G4 of
Definition~\ref{def:gap}). RQ2's answer: the Bridge Contract is an
execution-side precondition for GAR in non-GEC deployments. RQ3's
answer: the Bridge Contract's resume record and re-attach protocol complement
PACE; PACE handles ``how to re-attach'', the Bridge Contract handles
``what happened during the detach''.

\paragraph{Why a separate contract and not a profile or extension.}
An extension of one of AGTP, GAR or PACE would bind only the implementers
of that draft, yet each obligation of the contract crosses the scopes of
more than one of them: the transport (AGTP), the audit record (GAR) and
the resumption of a suspended session (PACE). A bridge sits at their
intersection, and none of AGTP, GAR and PACE can impose a requirement on the
others. Table~\ref{tab:why-separate} gives, for the obligations that the
matrix does not mark as specified in any source, and for audit continuity
where no GEC is present (GAR and the Audit Trail specify it for a GEC or
for a recorder other than the bridge), the closest provision as read and
why placing the obligation there would not suffice. Quarantine and restart
recovery are in the last row.

\begin{table}[htbp]
\centering
\caption{Why an existing draft, as read, does not suffice as the home of
an obligation.}
\label{tab:why-separate}
\footnotesize
\begin{tabularx}{\columnwidth}{p{0.24\columnwidth} X X}
\toprule
\textbf{Obligation} & \textbf{Closest provision} & \textbf{Why a profile of it is not enough} \\
\midrule
Detached execution & PACE defines a suspended state and recovery
(\S\ref{sec:rw-pace}); AGTP defines transport headers
(\S\ref{sec:rw-agtp}). & Neither says what the bridge does while the
client is away; the obligation is about a component that neither draft
defines. \\
Resume proof & PACE re-attaches a suspended session
but, as read, defines no resume credential and reserves
\texttt{SESSION\_RESUME} (0x06) without a body; the AGTP SUSPEND method as read issues
a single-use nonce for a requested pause but binds it to no audit
position, and the RESUME method itself is reported to be specified in the companion
\cite{I-D.hood-agtp-standard-methods} \S 8.7 (not read; reported by its table of contents) (\S\ref{sec:rw-pace}, \S\ref{sec:rw-agtp}). & The
proof depends on a durable, hash-chained record kept by the bridge
(I4, I5, I11, I12), not on a message field, so a message-level profile
would leave the record unspecified. \\
Audit continuity (no GEC) & GAR defines the record shape and presupposes a
GEC as the source (\S\ref{sec:rw-gar}). & A GAR profile cannot make the
bridge the source without dropping GAR's own assumption; the contract
instead emits into GAR's shape when no GEC is present
(Definition~\ref{def:gecless}). \\
Human-pending queue, authority drift &
HEM specifies a pending state, timeouts and single-use decisions per
governed object (\S\ref{sec:rw-hem}); none of the drafts ties them to held
tool calls. & The obligations couple the
escalation, the queue and the log (I2, I3, I13, I14); each draft would
carry only the part inside its own scope. \\
Quarantine, restart recovery & GAR quarantines telemetry spans from
unauthenticated processes (\S15), and the Audit Trail flags a chain break
and records a recovery gap (\S13.4, \S8.2, \S8.3); PACE and AGTP say that a
session survives a restart (\S\ref{sec:rw-gar}, \S\ref{sec:rw-pace}). &
None isolates a session whose log fails verification or re-verifies the
chain on restart; the obligations are about the bridge's own log, which a
profile of a message format would not reach. \\
\bottomrule
\end{tabularx}
\end{table}

This is an argument from the drafts as read. It would weaken if one of
them were extended to define the bridge component and its detach-window
obligations itself.

\paragraph{RQ4: contribution.} The contribution is one artifact made of a
specification, evidence and a gap analysis. The \emph{specification} is the naming of the orphaned bridge
obligations (G1--G4), the product state model, a normative transition
table, a challenge-response re-attach pinned to the acknowledged prefix of
the chain, and invariants I1--I15. The \emph{executable evidence} and the
\emph{formal model} are described in Section~\ref{sec:conformance} and
Appendix~\ref{app:formal}. Relative to
product implementations such as Herdr, which restore sessions but, as far as I could find, publish
no interoperable format or conformance criteria, the difference is that the Bridge
Contract is a specification that can be tested; it makes no claim that
such products lack the behaviour.

\begin{figure}[t]
\centering
\begin{tikzpicture}[x=1cm,y=1cm,font=\scriptsize,
  msg/.style={-{Stealth[length=4pt]},thick},
  note/.style={draw,rounded corners=1pt,fill=gray!10,inner sep=2pt,align=left},
  cell/.style={draw,rounded corners=1pt,fill=yellow!15,inner sep=2pt}]
\def\cx{0}\def\bx{3.1}\def\rx{6.2}\def\hx{8.9}
\foreach \n/\x in {Client/\cx,Bridge/\bx,Runtime/\rx,Human/\hx}{
  \node[draw,fill=blue!6,minimum width=1.6cm] (\n) at (\x,0) {\n};
  \draw[dashed,gray] (\x,-0.3) -- (\x,-11.4);}
\draw[msg] (\rx,-1) -- node[above]{escalate} (\bx,-1);
\draw[msg] (\bx,-1.6) -- node[above]{ask} (\hx,-1.6);
\node[cell,anchor=west] at (\bx+0.15,-2.2) {\texttt{HUMAN\_PENDING}: (Att,Pend)};
\draw[thick] (\cx-0.15,-3) -- (\cx+0.15,-3.3) (\cx-0.15,-3.3) -- (\cx+0.15,-3) ;
\node[anchor=west] at (\cx+0.3,-3.15) {transport lost};
\node[cell,anchor=west] at (\bx+0.15,-3.8) {\texttt{DETACHMENT}, \texttt{RESUME\_PREPARED}: (Det,Pend)};
\draw[msg] (\rx,-4.6) -- node[above]{tool call (queued, I2)} (\bx,-4.6);
\draw[msg] (\hx,-5.3) -- node[above]{answer; resolve} (\bx,-5.3);
\node[cell,anchor=west] at (\bx+0.15,-5.9) {\texttt{HUMAN\_RESOLVED}: (Det,None); no drain while away};
\draw[msg] (\cx,-6.7) -- node[above]{challenge} (\bx,-6.7);
\draw[msg] (\bx,-7.2) -- node[above]{nonce} (\cx,-7.2);
\draw[msg] (\cx,-7.9) -- node[above]{signed (nonce, seq, hash)} (\bx,-7.9);
\node[cell,anchor=west] at (\bx+0.15,-8.5) {\texttt{RESUME\_STARTED}: (Res,None); chain verified};
\node[cell,anchor=west] at (\bx+0.15,-9.2) {\texttt{RESUME\_VERIFIED}: (Att,None)};
\draw[msg] (\bx,-9.8) -- node[above]{events after seq} (\cx,-9.8);
\draw[msg] (\bx,-10.5) -- node[above]{drain: forward (I13, I14)} (\rx,-10.5);
\end{tikzpicture}
\caption{The disconnect / resume / re-attach sequence, with an escalation
pending when the client disconnects. Cells are written (attachment,
human). The re-attach is a challenge and a signed request, not a token;
the human's resolution while the client is away does not drain the queue,
verification does.}\label{fig:sequence}
\end{figure}

\subsection{Threat model (STRIDE-lite)}
\label{sec:threat-model}

This subsection presents the threat model for the Bridge Contract
following the STRIDE methodology~\cite{shostack2014threat}, restricted to the
categories relevant to bridge-mediated session governance (the
``-lite'' suffix). For each STRIDE category we state whether the
Bridge Contract mitigates it, partially mitigates it, or defers to
another specification; for ``No'' rows we identify the residual-risk
owner.

\paragraph{Assets.} The Bridge Contract protects the following
assets:
\begin{itemize}[leftmargin=2em]
\item \textbf{Event log integrity.} The append-only, hash-chained
log of governance events per session (A1, I5).
\item \textbf{Resume integrity and authenticity.} The resume record
written at detach and the challenge-response that gates re-attach (I1,
I4, I5, I9, I11).
\item \textbf{Authority scope.} The set of permissions granted to
the session at attach; widening is refused, or reported when the runtime does it itself, in every cell (I3, I10).
\item \textbf{Tool-call queue.} The queue of pending tool calls
while the human dimension is pending; queue contents are part of the resume
record and of the replay (I2, I13).
\item \textbf{Session lifecycle.} The set of allowed state
transitions; $\bot$ is terminal (I6).
\end{itemize}

\paragraph{Adversaries.} The threat model considers the following
adversaries:
\begin{itemize}[leftmargin=2em]
\item \textbf{Network attacker.} Can observe and modify traffic
between client, bridge, and (in GEC-present deployments) GEC.
The contract specifies no transport binding, so this is mitigated only by a transport the deployment supplies (``Approvals and resolutions are not signed'', Table~\ref{tab:limitations}); I-SIGN-EVENT is a SHOULD without a test.
\item \textbf{Byzantine client (in scope, \S\ref{sec:scope}).} May attempt to replay captured
re-attach requests, present a history that is not a prefix of the chain,
forge requests to terminate sessions, or supply malicious JWKS
endpoints. Mitigated by I4, I9 and I11 (prefix consistency, short-lived
single-use nonce, proof of possession, rejection without state change)
and I-SIGN-JWKS (MUST NOT accept session-supplied JWKS), assuming no
registered client key is stolen (A7).
The formal model gives this adversary explicit capabilities and checks
the re-attach path against it (Appendix~\ref{app:attacker}).
\item \textbf{Byzantine bridge.} May suppress governance events or
forge events. \emph{Out of scope} (A4); mitigations are
deployment-level (external transparency logs, e.g., SCITT), and the residual risk is owned by a future revision of the contract (see the register below).
\item \textbf{Operator error.} May misconfigure timeouts, JWKS
endpoints, or authority scope. Mitigated by I-SIGN-JWKS (JWKS
source restriction), I3 (authority lock) and the timer bounds of
Section~\ref{sec:timeouts}.
\item \textbf{Compromised runtime.} The agent runtime may be
prompt-injected or otherwise compromised. It may self-approve
escalations (bypassing I2/I13 if the bridge does not authenticate
the human channel), flood the tool-call queue (mitigated by
$Q_{\max}$, I13), squat on client keys (mitigated by I11-REVOKE
revocation), or attempt to widen its own authority (reported by I3
as \texttt{observed}, not prevented). The runtime is trusted for
execution (the bridge cannot see what the runtime does with a
forwarded call), but not for authority changes: I3 and I10 ensure
the bridge's record never widens, even if the runtime's effective
authority does.
\end{itemize}

\paragraph{Trust boundaries.} The trust boundaries assumed by the
Bridge Contract are:
\begin{itemize}[leftmargin=2em]
\item \textbf{Client $\leftrightarrow$ Bridge.} Untrusted: the
client may be Byzantine, but it does not control the bridge's event log;
its misbehaviour is mitigated by I4, I9, I11 and I-SIGN-JWKS, except where it
holds a valid key (A7).
\item \textbf{Bridge $\leftrightarrow$ Agent runtime.} Trusted
under A4 (the bridge is non-Byzantine); the agent runtime is
out of scope.
\item \textbf{Bridge $\leftrightarrow$ GEC.} In GEC-present
deployments, the GEC is the sole authoritative event source
(per GAR \S 3); the bridge forwards unsigned events.
\item \textbf{Bridge $\leftrightarrow$ Audit consumer.} The audit
consumer (e.g., a SIEM feed, a SCITT transparency log) reads
events emitted by the bridge; in GEC-less deployments, the
consumer verifies the operator key binding out-of-band
(I-SIGN-EVENT).
\end{itemize}

\paragraph{STRIDE coverage.} Table~\ref{tab:threats} maps STRIDE
categories to Bridge Contract invariants. ``S'' = Spoofing;
``T'' = Tampering; ``R'' = Repudiation; ``I'' = Information
Disclosure; ``D'' = Denial of Service; ``E'' = Elevation of
Privilege.

\begin{table}[htbp]
\centering
\caption{STRIDE coverage: threats the Bridge Contract mitigates
(or does not), with residual-risk owner for unmitigated rows.}
\label{tab:threats}
\small
\begin{tabularx}{\columnwidth}{l >{\raggedright\arraybackslash}X >{\raggedright\arraybackslash}X >{\raggedright\arraybackslash}X}
\toprule
\textbf{STRIDE} & \textbf{Threat} & \textbf{Mitigated by BC?} & \textbf{Residual-risk owner} \\
\midrule
S & Authority spoofing (forged or widened scope) & Partial (I3 observed, not enforced, when the runtime widens its own authority; a forwarded call in (Att,None) is not scope-checked; I10 for the bridge's record) & Runtime; deployment operator (Table~\ref{tab:limitations}) \\
S & JWKS endpoint spoofing & Yes (I-SIGN-JWKS, MUST; by inspection only, Table~\ref{tab:limitations})\textsuperscript{a} & --- \\
S & Re-attach spoofing & Yes (I11 proof of possession) & Key custody (A7): deployment \\
T & Event log tampering (edit or reordering of events inside the log) & Yes (A1, I5) & --- \\
T & Truncation of the log, or rewriting of its last events & No without a head held outside the log (I-CHECKPOINT anchor) & Bridge operator (storage integrity, external head) \\
T & Tampering with the history presented at re-attach & Yes (I4 prefix check, I5) & --- \\
T & Forged or injected approval, resolution or escalation on an unauthenticated transport (these ops carry no signature) & No: they are transport-authenticated (Appendix~\ref{app:wire-format}) & Transport binding; deployment operator \\
T & Tampering during storage write & Partial (I12, I5) & Bridge operator (storage integrity, A1) \\
T & Replay of queued side effects after detach or crash & Yes (I13), with a cooperating runtime & Runtime must honour idempotency keys \\
R & Agent repudiation & Partial (I5 + I-SIGN-EVENT) & GEC if present; otherwise bridge operator \\
R & Bridge repudiation (a Byzantine bridge is out of scope, A4) & Partial (I5 + I-SIGN-EVENT + I-CHECKPOINT; both SHOULD, not in the baseline; the reference bridges keep no external head and the Rust crate emits no checkpoints) & External transparency log (SCITT) \\
\textbf{I} & \textbf{Disclosure of the resume record and of the replay} & \textbf{Partial} & \textbf{Bridge operator (encryption at rest and in transit per \S\ref{sec:limitations})} \\
I & Event log information disclosure & No & Bridge operator (deployment-level access control) \\
I & Tool-call queue information disclosure & No & Bridge operator (queue contents in the log and the replay; mitigate via encryption) \\
D & Detach timeout exhaustion & Partial (L1 bounds detachment and gives no availability; Table~\ref{tab:limitations}) & Bridge operator; deployment (Table~\ref{tab:limitations}) \\
D & Human timeout exhaustion & Yes (L2) & --- \\
D & Session termination by forged re-attach attempts or by holding the nonce book & Partial (I11: a rejected request changes no cell; but a caller without a key that holds the live-nonce limit refuses a detached client every challenge until $T_{\text{detach}}$ ends the session, Section~\ref{sec:invariants-token-auth}) & Rate-limiting or authentication of the transport in front of \texttt{challenge}: bridge operator \\
D & Denying a genuine re-attach by flooding \texttt{challenge} or the log without a key & Partial (I11 (f), (g): a live nonce is never evicted, and the failures before authentication recorded between two verifying requests are bounded; a replayed authentic request is recorded and restores that allowance, so the log grows with the caller's effort; the limit of live nonces can still be held, delaying a new challenge by up to $T_{\text{nonce}}$ each time) & Rate-limiting or authentication of the transport in front of \texttt{challenge}: bridge operator \\
D & Bridge restart & Partial (I12, Section~\ref{sec:recovery}) & Bridge operator (HA, replication) \\
E & Authority drift in any cell & Partial (I3 observed; I10 for the bridge's record) & Runtime; deployment operator (Table~\ref{tab:limitations}) \\
E & Privilege escalation by presenting a stale or forged history & Yes (I4, I11) & --- \\
S & Replay of a captured re-attach or join request & Yes (a captured request cannot be replayed: I11 single use, I9 lifetime, and a nonce issued for a key that was not registered never verifies) & --- \\
S & Theft of a client key & Partial (I11-REVOKE: operator can revoke) & Key management: deployment (operator must sign revocation) \\
\bottomrule
\end{tabularx}
\end{table}

\noindent\textsuperscript{a} I-SIGN-JWKS is vacuous in GEC-present deployments (the GEC provides the JWKS endpoint).

\paragraph{Residual-risk register.} Rows marked ``Partial'' or
``No'' have residual risk owned by:
\begin{itemize}[leftmargin=2em]
\item \textbf{Bridge operator (deployment).} Configures
persistence, HA, encryption at rest, access control. Not in scope
for BC.
\item \textbf{GEC (if present).} Authoritative event source per
GAR; bridge forwards events in GEC-present deployments.
\item \textbf{External transparency log (SCITT).} For
non-repudiation of bridge events in GEC-less deployments where
the bridge operator is potentially untrusted.
\item \textbf{Future revision of BC.} For concerns explicitly
deferred (multi-hop composition, Byzantine bridge, partial
detach), the residual risk is owned by a future revision of this
specification.
\end{itemize}

\paragraph{Out-of-scope threats.} The following STRIDE-relevant
threats are explicitly out of scope for this document (see
\S\ref{sec:scope}):
\begin{itemize}[leftmargin=2em]
\item Bridge replication and failover (D): deployment-level.
\item Network partition / partial detach (D, T): not modelled.
\item True concurrency on the same session (T, R): events are
serialised (A5); interleavings are covered by the product machine only.
\item Byzantine bridge (T, R, I): A4 assumes non-Byzantine bridge.
\item Multi-hop composition (S, T, E): BC $\circ$ BC not modelled.
\end{itemize}

% ================================================================================
